\newpage
\section{State of the Art}


% Design principle for Chapter 2: Progressive funnel from broad RAG landscape → energy measurement → adaptive computation → module-level routing gap.
% Reader takeaway: "I understand the RAG design space, how energy is measured in ML systems, and why module-level routing under energy constraints is unexplored."


% ============================================================================
% 2.1 — RAG: Evolution and Architecture
% ============================================================================
\subsection{RAG Architectures}
\label{sec:architecture}

\subsubsection{Core Components}
% --- 2.1.1 Core Pipeline ---

Retrieval-Augmented Generation (RAG) systems link language models with content from an 
external knowledge source at inference time~\cite{rag2020}.
Large language models are trained once on public data; 
adapting this knowledge after training remains challenging. 
RAG systems address this limitation by retrieving relevant content 
from an external knowledge base to contextualize the model when a 
query arrives.


\begin{figure}[h!]
    \centering
        \centering
        \includegraphics[width=0.9\textwidth]{chart-rag-detail2.png}
        \caption{Core RAG components and their interactions}
        \label{fig:flow-rag-baby}
\end{figure}

Figure~\ref{fig:flow-rag-baby} illustrates the core components and their interactions.
A minimal RAG pipeline is composed of three core components:
\begin{enumerate}
    \item An \emph{embedding model} that transforms an input query into a dense vector
    representation.
    \item A \emph{vector database} that stores indexed document representations and
    executes similarity search against the query vector, returning a ranked list of
    candidate passages.
    \item A \emph{large language model} that receives the query together with the retrieved
    passages and generates the final response.
\end{enumerate}

The embedding step translates text into a vector format that enables efficient 
similarity search against the database, which stores documents in the same vector 
space. Embedding models are typically language models optimized for semantic 
search tasks.

From a functional perspective, retrieval serves as a mechanism for grounding and 
constraining generation: it can improve relevance, factual accuracy and reduce
hallucination rates. At a high level, the execution flow follows a simple sequence:

\begin{figure}[h]
    \centering
        \centering
        \includegraphics[width=0.85\textwidth]{flow-rag-simple-bw.png}
        \caption{Simple RAG sequence-flow: Embedding $\rightarrow$ Retrieval 
        $\rightarrow$ Generation}
        \label{fig:flow-rag-simple}
\end{figure}

A RAG system is best understood as an integrated pipeline rather than a single model. 
Its behavior and performance result from the interaction of multiple components, each 
with distinct computational and latency costs. System-level design choices therefore 
directly affect both output quality and resource consumption.

This thesis focuses on the inference behavior of RAG systems, i.e., 
the operational phase in which the system processes user queries and generates answers. 
Training procedures, document ingestion and model adaptation are not actively studied.

% --- 2.1.2 Enhancement Modules in Advanced RAG ---

\subsubsection{Advanced Enhancements}

The core pipeline can produce an answer for any query. Since its 
introduction in 2020~\cite{rag2020}, subsequent research has proposed 
numerous extensions to improve retrieval and answer quality.
Structurally, enhancements fall into two categories:

\begin{itemize}
    \item \emph{Modify an existing stage}: for example, switching retrieval 
    from dense to hybrid changes how retrieval operates without adding a new pipeline step.
    \item \emph{Add a new stage}: for example, query expansion inserts 
    an additional model call before retrieval.
\end{itemize}

This distinction matters for energy analysis: modifying a stage changes its cost profile, 
while adding a stage 
introduces an entirely new cost that was previously absent.
% These improvements can three phases: query formulation, retrieval, and post-retrieval 
% refinement.

\paragraph{Query Enhancement (adds a stage before retrieval).}
These methods reformulate or enrich the query before it searches for context. 
HyDE~\cite{hyde2023} generates a hypothetical answer and embeds it 
as the query representation, bridging 
the vocabulary gap between short queries and stored documents. 
Related approaches include: 

\begin{itemize}
    \item Step-Back Prompting~\cite{stepback2023} abstracts the query into a 
    higher-level question about underlying principles before retrieval. 
    \item Query Rewriting~\cite{queryrewriting2023} reformulates the query 
    into a more retriever-friendly form.
    \item Multi-query approaches generate several reformulations and 
    fuse their retrieval results via Reciprocal Rank Fusion.
    \item For multi-hop questions, query decomposition frameworks such as 
    IRCoT~\cite{ircot2022} interleave retrieval with chain-of-thought
    reasoning, breaking complex queries into sequential sub-queries.
\end{itemize}

All query enhancement methods share a cost property: they add at least 
one model call before retrieval begins.


\paragraph{Retrieval Strategy (modifies the retrieval stage).}
Rather than adding stages, retrieval improvements change \emph{how} the existing retrieval stage 
operates. Three main approaches exist:
\begin{itemize}
    \item \emph{Sparse retrieval} matches on keywords without embeddings. 
    (Classic: BM25~\cite{bm252009}, Modern: SPLADE~\cite{splade2021})
    \item \emph{Dense retrieval} embeds queries and documents into a shared 
    vector space and retrieves by similarity. Advances include late-interaction 
    models (ColBERTv2~\cite{colbertv2}) that compute token-level similarity 
    for higher accuracy and variable-dimension embeddings 
    (Matryoshka~\cite{matryoshka2022}) that allow coarse-to-fine search.
    \item \emph{Hybrid retrieval} combines dense and sparse signals 
    via score fusion, capturing both semantic and lexical relevance. 
    Reciprocal Rank Fusion (RRF)~\cite{rrf2009} is the standard 
    merging strategy. RAG-Fusion~\cite{ragfusion2023} extends this 
    idea by generating four to five query variations and fusing their 
    retrieval results into a single ranked list.
\end{itemize}
% At the index level, vector quantization techniques such as RaBitQ~\cite{rabitq2024} compress stored 
% embeddings by up to 32$\times$ with bounded recall loss. These modifications change the retrieval cost 
% profile (e.g.\ hybrid is slightly more expensive than dense-only) but do not add new pipeline stages.

\paragraph{Post-Retrieval Refinement (adds or modifies stages after retrieval.)}
After retrieval, additional processing can improve the context provided to generation. 
\begin{itemize}
    \item \emph{Reranking} uses a cross-encoder~\cite{nogueira2019} to reorder 
    candidates by relevance, keeping only the top-$n$ passages.
    \item \emph{Context compression}~\cite{llmlingua2023} removes low-information 
    tokens from retrieved passages, reducing prompt length and noise.
    \item \emph{Strategic reordering}~\cite{lostmiddle2023} rearranges 
    passages to place the most relevant at the beginning and end of the 
    prompt, counteracting the ``lost-in-the-middle'' degradation observed in LLMs.
    \item \emph{Post-retrieval verification} (Self-RAG~\cite{selfrag2024}) 
    generates reflection tokens to assess whether retrieved context is relevant 
    and supported, discarding poor chunks before generation. 
\end{itemize}   

Danilevsky et al.~\cite{llmintrinsics2025} implement individual RAG functions ; reranking, 
hallucination detection, citation generation, as separate LoRA adapters that can be composed or omitted per 
query, illustrating that these functions can be composed independently.

\medskip

These enhancements share a structural property: each can be independently enabled or disabled; when implemented
with this logic. Each has a measurable cost (latency, energy) and a conditional
benefit that depends on the query and context. 
Whether a technique adds a new stage or modifies an existing one, the cost is incurred for every query when 
the module is enabled (regardless of whether that specific query benefits from it). This modular, additive 
structure is central to the design space explored in this thesis.


% --- 2.1.3 From Static to Dynamic ---

\subsubsection{Modular RAG}

In most deployments, module activation is defined through static configuration
at deployment time. The same set of modules is applied to all queries regardless
of their complexity or expected benefit. 
This ``always-on'' design favors simplicity and predictability. 
Usually justified by a quality-first perspective in which the cost 
of executing all modules is assumed acceptable as long as answer quality improves.

Recent work illustrates this trend: LLM Intrinsics~\cite{llmintrinsics2025} demonstrates 
composable LoRA adapters for individual RAG functions and knowledge graph-enhanced 
systems like LightRAG~\cite{lightrag2025} introduce additional optional retrieval 
strategies. Recent commercial systems have also adopted automatic routing mechanisms
that select among models based on inferred query complexity, indicating growing
interest in adaptive execution beyond research settings.

Existing work on RAG systems focuses primarily on retrieval quality and answer 
relevance. Efficiency considerations, when present, are typically addressed indirectly 
through faster models or reduced model size. The cost of always-on module activation 
is rarely questioned: it is treated as a deployment-time decision, not a runtime 
optimization variable. The following section reviews how energy consumption is 
measured and analyzed in machine learning systems.



% ============================================================================
% 2.2 — The Energy Anatomy of RAG
% ============================================================================
\subsection{Energy in ML Systems}
\label{subsec:energy}

% Design principle for 2.2: Block A = energy as RAG-specific concern; Block B = stage-level cost asymmetry; Block C = measurement-to-control bridge.
% Reader takeaway: "RAG stages span orders of magnitude in energy cost. Component-level measurement makes each stage a controllable variable."

% --- Block A: Energy as a RAG-specific concern ---

Energy consumption is a distinct metric from latency or throughput. A system
can be fast yet energy-inefficient or slow but power-frugal. For organizations
running continuous inference workloads, operational costs scale directly with power
consumption~\cite{energyml2024}. On-premise deployments face hard power budgets
and thermal constraints that latency metrics alone do not capture.

RAG systems have a heterogeneous energy profile that single-model benchmarks miss.
Data movement across retrieval nodes and computation on GPU-bound generation run on
different hardware with different power characteristics. Standard inference metrics
(tokens per second, time to first token) measure the language model in isolation.
They do not capture the distributed cost of embedding, retrieval or orchestration.
Recent analysis confirms this: in RAG-based systems with agentic control flow, CPU-bound
operations account for 90\% of end-to-end latency and 44\% of total energy
consumption~\cite{cpucentric2025}. This distribution differs from what single-model benchmarks would suggest.

% --- Block B: The asymmetry of RAG stages ---

Not all pipeline stages carry equal cost. A keyword lookup (BM25) consumes millijoules
on a CPU. Generating a hypothetical answer for query expansion (HyDE) requires a full
LLM inference pass, consuming orders of magnitude more energy on a
GPU~\cite{rago2025}. Embedding a query takes milliseconds; a cross-encoder reranking
pass adds moderate cost; switching from dense to hybrid retrieval adds near-zero marginal
energy. 
% Chapter~\ref{sec:baseline_analysis} quantifies these per-stage costs.

% This asymmetry is structural: it follows from combining cheap information retrieval
% operations with expensive neural generation within a single request. 
Applying all stages uniformly to every query drives unnecessary cost when some
stages do not improve the answer for that query. To exploit this cost asymmetry, 
cost must first be measurable per stage.

% --- Block C: Measurement-to-control bridge ---

\begin{figure}[h!]
    \centering
        \centering
        \includegraphics[width=0.65\textwidth]{energy-levels.png}
        \caption{Energy measurement levels: system, component, and software}
        \label{fig:energy-levels}
\end{figure}

Energy measurement occurs at three levels:
\begin{itemize}
    \item \textbf{System-level} captures total power at the wall outlet; ground truth,
    but no attribution to individual components or operations.

    \item \textbf{Component-level} uses hardware counters: Intel RAPL provides CPU and
    DRAM energy; NVIDIA NVML exposes GPU power via \texttt{nvidia-smi}. Frameworks such as
    MELODI~\cite{melodi2025} and APEX+~\cite{apexplus2025} build on these interfaces.
    Thamm and Leser~\cite{raplenergy2025} evaluate RAPL measurement strategies in
    distributed environments. This is directly relevant to multi-node RAG deployments.

    \item \textbf{Software-level} estimates approximate energy from proxy metrics (FLOPs,
    memory accesses, execution time) but introduce modeling error.
\end{itemize}

Component-level measurement provides an appropriate granularity for RAG systems. It attributes energy
to the CPU package and GPU device on each participating node. In a RAG pipeline, this
means energy can be attributed to query expansion, retrieval, reranking and generation
separately. If cost is measurable per stage, each stage becomes a controllable variable.
This enables component-level optimizations.

Guo et al.~\cite{keplerenergy2026} conduct a controlled experiment on energy reduction
in RAG systems, confirming that configuration choices measurably affect total consumption.
Their work evaluates fixed technique combinations rather than per-query adaptation. Single-model optimizations such as quantization and attention kernel
optimization~\cite{flashattn2022} reduce energy within one serving framework.
T. Burghardt and S. Chou~\cite{burghardt2025} applied quantization and fine-tuning to the
generation model of the same knowledge management system used in this thesis, reducing
memory by 39\% and inference time by 57\%; targeting the most energy-intensive RAG
stage. However, these techniques are model-specific: each new model generation requires
repeating the compression and tuning pipeline.


% ============================================================================
% 2.3 — Adaptive Orchestration and Module Routing
% ============================================================================
\subsection{Adaptive Orchestration}
\label{subsec:adaptive}

% Design principle for 2.3: Block A = conditional computation paradigms; Block B = model routing ("static prompt" limitation); Block C = adaptive RAG (quality-first limitation); Block D = module-level routing definition; Block E = research gap.
% Reader takeaway: "Conditional execution is proven. Model routers ignore pre-generation cost. Pipeline routers optimize quality, not energy. Module-level routing under energy constraints is unexplored."


% --- 2.3.1 Conditional Computation Paradigms ---

\subsubsection{Conditional Computation}
\label{subsec:conditional}

Conditional computation addresses the inefficiency of the full computational pipeline for 
every input by selectively executing 
parts of a system. It does so based on input characteristics or intermediate signals~\cite{moe2022}. 
The challenge lies in deciding \emph{when} to activate additional computation and 
\emph{when} to skip it, without degrading output quality.

Several forms of adaptive execution have been explored:
\begin{itemize}
    \item \textbf{Early exit mechanisms} allow deep networks to terminate inference at intermediate layers once sufficient confidence is reached, avoiding computation in deeper layers~\cite{branchynet2016, deebert2020}.
    \item \textbf{Mixture-of-experts (MoE)} models route each input to a subset of specialized submodels, activating only part of the network per input~\cite{moe2022}.
    \item \textbf{Cascaded systems} process inputs using a sequence of models with increasing complexity. Inexpensive models handle easy cases; only difficult inputs are forwarded to costly ones. Dekoninck et al.~\cite{unifiedrouting2025} formalize both routing and cascading as instances of linear cost-quality optimization and prove that cascade policies dominate fixed-threshold cascading.
\end{itemize}

These mechanisms share a common structure: a lightweight decision function gates 
access to expensive computation. Decision logic is typically implemented using heuristics, 
confidence thresholds or lightweight classifiers. Evaluation in prior work primarily 
targets accuracy--latency trade-offs, with energy efficiency treated indirectly through 
faster inference or reduced computation. Adaptation is usually applied within 
a single model or architecture rather than across multiple system components.


% --- 2.3.2 LLM Routing: The "Static Prompt" Limitation ---

\subsubsection{Model Routing}

A direct application of conditional computation to language models is \emph{model routing}: 
given a query, select which LLM should answer it. The query remains unchanged; only the 
answering model varies. Theoretical analysis shows that an oracle router, one that perfectly 
predicts model performance, can match the quality of the best model in the pool while 
reducing inference cost by 50--85\%~\cite{routerbench2024}.

Several frameworks operationalize this idea. RouteLLM~\cite{routellm2025} trains a 
Bradley-Terry preference model combined with matrix factorization to predict the win 
probability of a strong model over a weak one, using query embeddings as input. 
FrugalGPT~\cite{frugalgpt2023} implements cascading with agreement-based stopping: a 
cheap model answers first and its response is accepted if confidence exceeds a 
threshold; otherwise, the query escalates to a more expensive model, achieving up 
to 98\% cost reduction. Dekoninck et al.~\cite{unifiedrouting2025} provide theoretical 
grounding by formalizing model selection as linear cost-quality optimization.

The predictor signals used in model routing span multiple levels of complexity. 
Dense query embeddings capture semantic domain and topic~\cite{routellm2025}. Explicit 
features (token count, constraint density, code markers, question type) provide 
interpretable complexity indicators. Uncertainty-based methods use model confidence as 
a routing signal: Zhang et al.~\cite{uncertaintyrouting2025} route queries based on 
semantic entropy, a measure that clusters sampled responses by meaning before computing 
uncertainty and report 4--8\% accuracy gains over preference-trained routers on RouterBench. 
Predictor quality bounds routing quality: a router can only be as good as its ability 
to estimate task difficulty before execution.

These routers share a structural assumption: the prompt is already prepared. 
They select which model generates the answer but do not account for the energy 
spent \emph{constructing} that prompt, i.e., the Joules consumed by 
query expansion, retrieval or reranking before generation begins.
Usually, they optimize generation cost, not pipeline cost.


% --- 2.3.3 Adaptive RAG: The State of the Art ---

\subsubsection{Adaptive RAG}
\label{subsec:adaptive-rag}

Model routing selects \emph{which} model answers a query. A different line of 
work adapts \emph{how} the query is processed within the pipeline; deciding whether 
to retrieve, how to retrieve and what post-processing to apply. The pipeline itself 
becomes the design space.

Adaptive-RAG~\cite{adaptiverag2024} classifies queries by complexity (simple, medium, 
complex) and routes them to different retrieval strategies accordingly. 
Self-RAG~\cite{selfrag2024} introduces special reflection tokens (Retrieve, 
IsRelevant, IsSupported, IsUseful) that the model generates during inference to 
decide at each step whether retrieval is needed and whether retrieved context is 
relevant. This enables fine-grained control but introduces token overhead. CRAG 
(Corrective RAG)~\cite{crag2024} evaluates retrieved document relevance with a 
lightweight evaluator and triggers web search as a fallback when confidence is 
low. This is a strategy that can cause energy spikes when the fallback is invoked.

At the routing layer, RAGRouter~\cite{ragrouter2025} uses contrastive learning to 
match queries with the LLM configuration that produces the best answer given retrieved 
documents, achieving 3.29\% quality gains over static routing. 
UniRoute~\cite{uniroute2025} takes a training-free approach, clustering queries 
and routing based on cluster-level performance statistics, enabling zero-shot 
generalization to unseen models. DAAO~\cite{daao2025} applies difficulty-aware 
orchestration to agent workflows, achieving 6$\times$ cost reduction by matching 
queries to appropriately-sized models.

These systems are quality-first routers: when uncertain, they escalate computation 
to maximize answer quality. CRAG exemplifies this pattern: uncertainty triggers web 
search, causing an energy spike with no budget awareness. None use measured energy 
(Joules) as the optimization objective. Module activation (whether HyDE should be 
enabled, whether reranking is worth its cost for a specific query) is not treated as 
the decision variable. These systems escalate computation to maximize quality.
The present work instead constrains computation under an energy budget.


% --- 2.3.4 Defining Module-Level Routing ---

\subsubsection{Module-Level Routing}

This thesis defines module-level routing: not model selection (choosing between LLMs),
not retrieval gating (deciding whether to retrieve), but \emph{configuration 
vector control}; selecting the intensity of the pipeline per query. The decision 
variable is a configuration vector: \{HyDE: on/off, reranking: on/off, search type: 
dense/hybrid\}. Each query receives multiple independent binary decisions, not a 
single model choice.

Module-level routing differs from existing approaches in three ways. First,
energy is the primary objective (minimize Joules), with quality as a constraint---not 
the reverse. Second, prediction must happen \emph{before} pipeline execution, using 
only pre-execution query features (text characteristics, NLP signals, semantic 
embeddings). Third, the router must be lightweight enough that its own cost is 
negligible relative to the stages it gates.


% --- 2.3.5 Research Gap ---

\subsection{Research Gap}

RAG systems are modular and expose optional modules that can be independently
toggled (Section~\ref{sec:architecture}). Energy is measurable per stage and stages vary by orders
of magnitude in cost (Section~\ref{subsec:energy}). Conditional computation and query routing
are proven paradigms for reducing unnecessary work, with mature implementations
for model selection and emerging ones for pipeline-level adaptation
(Sections~\ref{subsec:conditional}--\ref{subsec:adaptive-rag}).

Despite these advances, module-level activation in RAG pipelines under explicit 
energy constraints remains relatively unexplored. Most existing routers optimize accuracy, latency or 
cost-per-token but do not address the energy cost of optional pipeline stages. 
Existing energy-aware methods optimize individual models but do not cross component 
boundaries. Existing work typically addresses these aspects separately:
modular RAG design, component-level energy measurement and adaptive computation
are not often combined to optimize module activation per query under energy constraints. 
This gap motivates the system-level approach developed in this thesis.



% ============================================================================
% 2.4 — Quality Evaluation in RAG
% ============================================================================
\subsection{Quality Evaluation}
\label{subsec:quality-eval}

Evaluating RAG output quality is a multi-dimensional problem. Deterministic metrics 
such as Exact Match, Token F1 and ROUGE-L capture surface-level overlap between 
generated answers and reference answers but miss semantic equivalence. Two answers 
can convey the same meaning using different words. Retrieval-specific metrics 
(R-Precision, Recall@$k$) evaluate the retrieval stage in isolation without assessing 
end-to-end answer quality.

LLM-as-judge approaches address these limitations by using a language model to
evaluate answer quality along specific dimensions. RAGAS~\cite{ragas2023} provides
automated evaluation along four axes: faithfulness (is the answer grounded in retrieved
context?), factual correctness (does the answer match the reference?), context precision
and context recall. A recent survey of 582~papers confirms that traditional metrics
(F1, EM, ROUGE) still dominate RAG evaluation, but LLM-based methods are adopted at an
increasing rate~\cite{ragevalsurvey2025}.

Evaluator model choice affects reliability. Wang~et~al.~\cite{vrageval2024}
report that GPT-4 achieves 82.6\% agreement with human annotations on
binary accept/reject decisions, whereas an 8B-parameter model achieves
only 36.8\%. Smaller evaluators tend to compress scores toward a
single label, reducing discriminative power. The evaluator model must
also differ from the generator to avoid self-evaluation bias, where
models favour outputs that match their own generation
patterns~\cite{vrageval2024}.

The choice of evaluation framework, evaluator model and bias-mitigation
strategy used in this thesis is detailed in Chapter~\ref{sec:method_protocol}.